\documentclass[a4paper,12pt]{article}

\usepackage{fontspec}
\setmainfont{PT Serif}
\usepackage{babel}
\babelprovide[import,main]{russian}
\usepackage{microtype}
\usepackage[margin=2.5cm]{geometry}

\setlength{\parindent}{1.25cm}
\setlength{\parskip}{0pt}

\begin{document}

\begin{center}
    {\Large\bfseries Теория вероятностей}

    \vspace{0.3em}

    {\large Семинар 5}

    \vspace{0.3em}

    \textbf{Условная вероятность}

    \vspace{0.3em}

    {\small 12 сентября 2026 года}
\end{center}

\vspace{1em}

\section*{Задача 1}

Условие первой задачи.

\textbf{Решение.}

Решение первой задачи.

\section*{Краткое напоминание}

Небольшой теоретический блок расположен там, где он встречается в исходном
материале.

\section*{Задача 2}

Условие второй задачи.

\textbf{Решение.}

Решение второй задачи.

\medskip

\noindent\textbf{Замечание.}
Реплика преподавателя сохраняется как отдельный блок.

\section*{Задача 3}

Условие третьей задачи.

% Конкретное оформление задачи, решения и замечания определяется bricks.

\end{document}
